% --- INSTITUTIONAL TEMPLATE FED v7.4.9 ---
% Estándar de Reporteo APA 7 (Babel-Free)
\documentclass[12pt,spanish]{article}

\usepackage{graphicx}
\usepackage{geometry}
\usepackage{indentfirst} 
\usepackage{caption}
\usepackage{floatrow} 

% Forzar estilos de caption antes de cualquier carga
\floatsetup[figure]{capposition=top}
\floatsetup[table]{capposition=top}

\usepackage{booktabs}
\usepackage{float}
\usepackage{longtable}
\usepackage{threeparttable}
\usepackage{array}
\usepackage{calc}
\usepackage{amsmath, amssymb}
\usepackage{fancyhdr}
\usepackage{mathptmx} 
\usepackage[T1]{fontenc}
\usepackage{setspace} 
\usepackage{ragged2e} 
\usepackage{titlesec}
\usepackage{url}
\def\UrlBreaks{\do\/\do-\do\.\do\_\do\?\do\=\do\&\do\%}
\usepackage[hidelinks,breaklinks]{hyperref}

% --- Pandoc Compatibility Block (Iron Clad) ---
\providecommand{\tightlist}{%
  \setlength{\itemsep}{0pt}\setlength{\parskip}{0pt}}

\usepackage{xcolor}
\usepackage{fancyvrb}
\usepackage{framed}

% --- Code Highlighting (Pandoc Standard) ---
\AtBeginDocument{
  \definecolor{shadecolor}{RGB}{248,248,248}
}

% Compatibilidad Pandoc 3.x
\ifdefined\pandocbounded\else
  \newcommand{\pandocbounded}[1]{#1}
\fi

% --- Pandoc Citeproc (CSL) compatibility ---
\newlength{\cslhangindent}
\setlength{\cslhangindent}{1.27cm} 
\newlength{\csllabelwidth}
\setlength{\csllabelwidth}{3em}
\newcommand{\citeproctext}{} 
\newcommand{\citeprocitem}[2]{\item} 
\newenvironment{CSLReferences}[2] 
 {\begin{list}{}{
  \setlength{\leftmargin}{\cslhangindent}
  \setlength{\parsep}{0pt}
  \setlength{\itemsep}{6pt}
  \ifodd #1
   \setlength{\leftmargin}{\leftmargin + \csllabelwidth}
   \setlength{\itemindent}{-\csllabelwidth}
  \fi
  }
  \renewcommand{\bibitem}[2][]{\item} % Elimina los corchetes [ ]
 }
 {\end{list}}
\newcommand{\CSLBlock}[1]{#1\hfill\break}
\newcommand{\CSLLeftMargin}[1]{\parbox[t]{\csllabelwidth}{#1}}
\newcommand{\CSLRightInline}[1]{\parbox[t]{\linewidth - \csllabelwidth}{#1}\break}
\newcommand{\CSLIndent}[1]{\hspace{\cslhangindent}#1}

\makeatletter
\@ifundefined{paragraph}{\newcounter{paragraph}}{}
\@ifundefined{subparagraph}{\newcounter{subparagraph}}{}
\makeatother

% Solución al error 'No counter none defined' en longtable
\makeatletter
\newcounter{none}
\@ifpackageloaded{caption}{
  \DeclareCaptionType{none}
  \captionsetup[none]{name=Tabla}
}{}
\makeatother

\hypersetup{
  colorlinks=true,
  linkcolor=blue,
  filecolor=blue,
  citecolor=blue,
  urlcolor=blue
}

% --- Macros Institucionales ---
\newcommand{\fuente}[1]{
  \vspace{4pt}
  \begin{flushleft}
    \small
    \textit{Nota.} #1
  \end{flushleft}
}

% --- Localización Manual Determinista (Sin Babel) ---
\AtBeginDocument{
  \renewcommand{\tablename}{Tabla}
  \renewcommand{\figurename}{Figura}
  \renewcommand{\contentsname}{Índice}
  \renewcommand{\listtablename}{Índice de Tablas}
  \renewcommand{\listfigurename}{Índice de Figuras}
  \renewcommand{\abstractname}{Resumen}
  \renewcommand{\refname}{Referencias}
  
  \RaggedRight
  \setlength{\parindent}{1.27cm}
  \setlength{\parskip}{0pt}
}

\DeclareCaptionLabelSeparator{newline}{\\ \vspace{2pt}}
\captionsetup{
  position=top,
  justification=RaggedRight,
  singlelinecheck=false,
  labelfont=bf,
  textfont=it,
  labelsep=newline,
  font=small,
  skip=10pt
}
\captionsetup[figure]{name=Figura}
\captionsetup[table]{name=Tabla}

% Márgenes APA 7
\geometry{
  a4paper,
  margin=25.4mm,
  headheight=15mm
}

% --- Encabezados ---
\pagestyle{fancy}
\fancyhf{} 
\fancyhead[L]{\includegraphics[height=1.2cm]{\logoUNL}}
\fancyhead[R]{\includegraphics[height=1.2cm]{\logoECONOMIA}}
\renewcommand{\headrulewidth}{0pt}
\fancyfoot[R]{\thepage}

\fancypagestyle{plain}{
  \fancyhf{}
  \fancyhead[L]{\includegraphics[height=1.2cm]{\logoUNL}}
  \fancyhead[R]{\includegraphics[height=1.2cm]{\logoECONOMIA}}
  \fancyfoot[R]{\thepage}
  \renewcommand{\headrulewidth}{0pt}
}

% --- Títulos ---
\titleformat{\section}{\normalfont\normalsize\bfseries\centering}{\thesection}{1em}{}
\titleformat{\subsection}{\normalfont\normalsize\bfseries}{\thesubsection}{1em}{}
\titleformat{\subsubsection}{\normalfont\normalsize\bfseries\itshape}{\thesubsubsection}{1em}{}
\titleformat{\paragraph}{\normalfont\normalsize\bfseries}{\theparagraph}{1em}{}
\titlespacing*{\section}{0pt}{12pt}{6pt}
\titlespacing*{\subsection}{0pt}{12pt}{6pt}
\titlespacing*{\subsubsection}{0pt}{12pt}{6pt}
\titlespacing*{\paragraph}{0pt}{12pt}{6pt}

\newenvironment{referenciasAPA}{
  \setlength{\parindent}{-1.27cm}
  \setlength{\leftskip}{1.27cm}
  \setlength{\parskip}{6pt}
}{\par}

\makeatletter\@ifpackageloaded{caption}{}{\usepackage{caption}}\AtBeginDocument{%
\ifdefined\contentsname
  \renewcommand*\contentsname{Table of contents}
\else
  \newcommand\contentsname{Table of contents}
\fi
\ifdefined\listfigurename
  \renewcommand*\listfigurename{List of Figures}
\else
  \newcommand\listfigurename{List of Figures}
\fi
\ifdefined\listtablename
  \renewcommand*\listtablename{List of Tables}
\else
  \newcommand\listtablename{List of Tables}
\fi
\ifdefined\figurename
  \renewcommand*\figurename{Figure}
\else
  \newcommand\figurename{Figure}
\fi
\ifdefined\tablename
  \renewcommand*\tablename{Table}
\else
  \newcommand\tablename{Table}
\fi
}\@ifpackageloaded{float}{}{\usepackage{float}}
\floatstyle{ruled}
\@ifundefined{c@chapter}{\newfloat{codelisting}{h}{lop}}{\newfloat{codelisting}{h}{lop}[chapter]}
\floatname{codelisting}{Listing}\newcommand*\listoflistings{\listof{codelisting}{List of Listings}}\makeatother\makeatletter\makeatother\makeatletter\@ifpackageloaded{caption}{}{\usepackage{caption}}
\@ifpackageloaded{subcaption}{}{\usepackage{subcaption}}\makeatother

\begin{document}

\begin{center}
    {\bfseries \Large UNIVERSIDAD NACIONAL DE LOJA}\\[0.3cm]
    {\bfseries \large FACULTAD JURÍDICA, SOCIAL Y
ADMINISTRATIVA}\\[0.2cm]
    {\bfseries CARRERA DE ECONOMÍA}\\[1.5cm]
    
    {\bfseries \large POLÍTICA ECONÓMICA}\\[1cm]
    
    {\bfseries \Large }\\[1.5cm]
\end{center}

\begin{center}
    \setstretch{1.1}
                        Erick Fabricio Condoy Seraquive \\
                \vspace{0.5cm}
    Econ. Maria Gabriela Moreno Hurtado \\ 
    
\end{center}

\vspace{1cm}

\doublespacing
\setlength{\parindent}{1.27cm}

\section{Guía de Transcripción a Mano: Explicación de la Política
Monetaria}\label{guuxeda-de-transcripciuxf3n-a-mano-explicaciuxf3n-de-la-poluxedtica-monetaria}

\begin{quote}
{[}!NOTE{]} Este documento contiene la síntesis académica rigurosa que
debes copiar a mano para cumplir con la justificación obligatoria de una
página. Se ha optimizado para que quepa físicamente en una página
escrita, manteniendo una estructura conceptual de alta densidad sin
redundancias ni ``slop''.
\end{quote}

\begin{center}\rule{0.5\linewidth}{0.5pt}\end{center}

\subsection{EXPLICACIÓN ESCRITA A MANO (Copiar esta
sección)}\label{explicaciuxf3n-escrita-a-mano-copiar-esta-secciuxf3n}

\textbf{Título:} Justificación Estructural y Análisis Teórico del
Capítulo 11: Política Monetaria

La estructura del organizador gráfico se fundamenta en un
\textbf{enfoque de transmisión causal e instrumental}. Este diseño
jerárquico permite trazar de forma lógica la secuencia de la
intervención pública desde el control técnico de variables operativas
hasta la consecución de los objetivos finales de estabilización
macroeconómica, separando las fases de implementación, transmisión y
retroalimentación (expectativas).

\textbf{1. Idea central del capítulo 11:} La idea central es que la
política monetaria, bajo el mando de una autoridad autónoma que genera
credibilidad y gestiona expectativas, busca estabilizar variables reales
(producción y empleo) y nominales (control de la inflación) mediante la
manipulación de la liquidez, el crédito y la tasa de interés en el
sistema financiero.

\textbf{2. Relación entre instrumentos, canales y objetivos:} Los
\textbf{instrumentos} (operaciones de mercado abierto, tasa de encaje y
redescuento) modifican las reservas bancarias primarias y el tipo de
interés interbancario. Esta alteración se propaga a través de los
\textbf{canales de transmisión} (tipo de interés, canal de crédito,
canal de precios de activos y tipo de cambio) impactando en el costo del
financiamiento y las decisiones de consumo e inversión de los agentes.
Finalmente, esta variación en la demanda agregada determina el nivel de
precios y empleo, alcanzando los \textbf{objetivos finales} de
estabilidad de precios e inflación controlada.

\textbf{3. Efectos distintos en el corto y largo plazo:} En el
\textbf{corto plazo}, debido a la existencia de rigideces nominales
(precios e información asimétrica), la política monetaria tiene efectos
reales sobre la producción y el empleo mediante alteraciones en el tipo
de interés real. En el \textbf{largo plazo}, bajo el supuesto de
neutralidad monetaria y flexibilidad de precios, los incrementos de la
oferta monetaria se trasladan de forma proporcional al nivel general de
precios, determinando únicamente la inflación sin alterar las variables
reales de la economía.

\textbf{4. Limitaciones del instrumento en una economía dolarizada
(Ecuador):} Ecuador carece de moneda propia y de un banco central emisor
de última instancia. Sus \textbf{limitaciones} clave son: (a)
incapacidad para emitir masa monetaria para inyectar liquidez directa,
(b) imposibilidad de manipular el tipo de cambio nominal para absorber
shocks externos (pérdida de canal externo de transmisión), y (c)
subordinación de la oferta monetaria a la balanza de pagos (ingreso de
divisas por exportaciones, inversión extranjera o deuda). La política
monetaria nacional queda reducida a medidas de liquidez secundaria
(encaje bancario y reciclaje interno del crédito), careciendo de
autonomía monetaria real.

\textbf{Conclusión (Respuesta al Cierre):} La política monetaria no
puede analizarse uniformemente entre países debido a las diferencias en
sus regímenes cambiarios y grados de desarrollo financiero. En regímenes
con soberanía monetaria se analiza bajo la libre fijación de tasas de
interés y tipo de cambio flexible, mientras que en Ecuador, al ser una
economía dolarizada, el análisis se desplaza hacia la gestión de la
liquidez existente, la solidez del sistema financiero y la
vulnerabilidad ante shocks de la cuenta corriente.

\begin{center}\rule{0.5\linewidth}{0.5pt}\end{center}

\subsection{Recomendaciones para tu Organizador
Gráfico}\label{recomendaciones-para-tu-organizador-gruxe1fico}

Para asegurar el \textbf{100\% de la rúbrica}, te sugiero estructurar tu
organizador (en Canva, draw.io o PowerPoint) de la siguiente manera:

\begin{enumerate}
\def\labelenumi{\arabic{enumi}.}
\tightlist
\item
  \textbf{Eje Central: La Ruta Causal} (Recomendado como organizador
  principal):

  \begin{itemize}
  \tightlist
  \item
    Dibuja un diagrama de flujo horizontal dividido en 4 niveles o
    fases:

    \begin{itemize}
    \tightlist
    \item
      \emph{Fase 1: Instrumentos} (OMA, Encaje Legal, Facilidades de
      Crédito).
    \item
      \emph{Fase 2: Variables Operativas e Intermedias} (Tasa
      interbancaria, Liquidez del sistema).
    \item
      \emph{Fase 3: Canales de Transmisión} (Crédito, Tasas de interés,
      Expectativas).
    \item
      \emph{Fase 4: Objetivos Finales} (Estabilidad de precios / Control
      de Inflación).
    \end{itemize}
  \end{itemize}
\item
  \textbf{Sección Lateral Izquierda (Fundamentos y Estrategias):}

  \begin{itemize}
  \tightlist
  \item
    \emph{Definición:} Acción sobre liquidez, crédito y tasas para
    influir en producción y precios.
  \item
    \emph{Autonomía del Banco Central:} Clave para la credibilidad de
    las reglas monetarias y evitar inconsistencia dinámica por presiones
    de financiamiento al gasto público.
  \item
    \emph{Estrategias:} Agregados monetarios frente a Metas de Inflación
    (\emph{Inflation Targeting}).
  \end{itemize}
\item
  \textbf{Sección Lateral Derecha (Efectividad y Límites):}

  \begin{itemize}
  \tightlist
  \item
    \emph{Efectividad:} Sensibilidad de la demanda de dinero (\(L(i)\))
    y de la inversión (\(I(r)\)).
  \item
    \emph{Limitaciones:} Retardos (interno y externo), asimetría (es más
    fácil frenar que reactivar), e inestabilidad de la velocidad del
    dinero.
  \end{itemize}
\item
  \textbf{Recuadro Especial (Ecuador - Dolarización):}

  \begin{itemize}
  \tightlist
  \item
    \emph{Título:} \textbf{¿Qué cambia en una economía dolarizada como
    Ecuador?}
  \item
    \emph{Puntos clave:}

    \begin{itemize}
    \tightlist
    \item
      No hay emisor de última instancia ni devaluación cambiaria.
    \item
      El Banco Central del Ecuador (BCE) no controla el tipo de interés
      internacional, sino que actúa como regulador de la liquidez
      doméstica (encaje mínimo y distribución de reservas).
    \item
      Dependencia del flujo neto de dólares externos.
    \end{itemize}
  \end{itemize}
\end{enumerate}

\end{document}
